
As versões de Heckbert e de Smelik et al.\ incluídas na pasta \texttt{docs/artigos/} são as versões dos autores (\textit{preprints}), cuja paginação difere da revista. As páginas citadas referem-se a essas versões, para que possam ser conferidas nos arquivos entregues. Os mesmos resumos estão em \texttt{docs/resumos/}.

\section{Heckbert (1986)}

\textbf{Título:} \textit{Survey of Texture Mapping}. \textbf{Autor:} Paul S. Heckbert. \textbf{Ano:} 1986 (IEEE Computer Graphics and Applications, v.~6, n.~11) \cite{heckbert1986}.

\textbf{Objetivo.} Organizar o conhecimento sobre texture mapping disponível até meados dos anos 1980, separando o problema em mapeamento geométrico e filtragem \cite[p.~1]{heckbert1986}.

\textbf{Principais conceitos.}
\begin{itemize}
    \item O mapeamento como composição de uma parametrização da superfície com a projeção para a tela \cite[p.~2]{heckbert1986}.
    \item Os parâmetros que podem ser mapeados (cor, normal, especular, transparência, deslocamento) \cite[p.~2]{heckbert1986}.
    \item Os mapeamentos afim, bilinear e projetivo, e a interpolação correta em perspectiva, com uma divisão por pixel \cite[pp.~4--5]{heckbert1986}.
    \item A taxonomia dos filtros (convolução direta, pré-filtragem e séries de Fourier) \cite[p.~7]{heckbert1986}, com uma tabela de custo e memória em que a pirâmide de mipmaps ocupa \(4/3\) da textura e a interpolação trilinear custa 8 acessos por pixel \cite[pp.~9--10]{heckbert1986}.
\end{itemize}

\textbf{Relação com o tema.} É a base do UV mapping: define o espaço \((u,v)\), mostra como ele chega a cada pixel e quanto custa ler a textura corretamente. O UV dinâmico altera justamente a parametrização descrita pelo autor. O artigo é anterior às GPUs programáveis e não trata de UV animado.

\textbf{Informações aproveitadas.} Composição do mapeamento e ordem de varredura (Cap.~\ref{chap:fundamentacao}); interpolação com correção de perspectiva e seleção de mipmap (Cap.~\ref{chap:equacoes}); aliasing, custo da filtragem trilinear e memória dos mipmaps (Caps.~\ref{chap:fundamentacao} e~\ref{chap:vulkan}).

\section{Dong et al. (2020)}

\textbf{Título:} \textit{Survey of Procedural Methods for Two-Dimensional Texture Generation}. \textbf{Autores:} Junyu Dong, Jun Liu, Kang Yao, Mike Chantler, Lin Qi, Hui Yu e Muwei Jian. \textbf{Ano:} 2020 (Sensors, v.~20, n.~4, 1135) \cite{dong2020}.

\textbf{Objetivo.} Revisar e comparar, pela aparência que produzem, os métodos de geração procedural de texturas bidimensionais estáticas. Texturas sólidas, jogos e animação ficam explicitamente fora do escopo \cite[p.~2]{dong2020}.

\textbf{Principais conceitos.}
\begin{itemize}
    \item A definição de textura procedural como imagem criada por um modelo matemático \cite[p.~1]{dong2020}.
    \item A classificação em métodos estruturados (autômatos celulares, reação-difusão, textons, Voronoi, texturas celulares), não estruturados (funções de ruído no domínio da frequência e no espacial, como Gabor e Perlin) e de pós-processamento \cite[pp.~3--4]{dong2020}.
    \item Uma avaliação perceptual de 20 métodos, com 450 amostras e 40 participantes \cite[pp.~19--20]{dong2020}, resumida em uma tabela de aplicações típicas de cada método \cite[p.~22]{dong2020}.
\end{itemize}

\textbf{Relação com o tema.} Fundamenta o lado procedural da técnica: o que é uma textura gerada por função, quando ela é vantajosa (tamanho ilimitado, sem costuras, famílias de variações) e quais são seus problemas (controle de parâmetros, aliasing, dificuldade de reproduzir texturas reais) \cite[pp.~2, 13, 23]{dong2020}.

\textbf{Informações aproveitadas.} Definição e taxonomia do procedural, critérios de quando usar e limitações (Cap.~\ref{chap:fundamentacao}); classificação dos nós procedurais do Blender (Cap.~\ref{chap:aplicacoes}); a ressalva de que o tempo relatado (menos de um segundo por imagem, em computador de mesa) mede a geração de uma imagem inteira, e não a avaliação por fragmento em tempo real \cite[p.~21]{dong2020}.

\section{Smelik et al. (2014)}

\textbf{Título:} \textit{A Survey on Procedural Modelling for Virtual Worlds}. \textbf{Autores:} Ruben M. Smelik, Tim Tutenel, Rafael Bidarra e Bedrich Benes. \textbf{Ano:} 2014 (Computer Graphics Forum, v.~33, n.~6) \cite{smelik2014}.

\textbf{Objetivo.} Revisar os métodos de modelagem procedural para mundos virtuais e entender por que eles ainda são pouco adotados por profissionais criativos \cite[p.~1]{smelik2014}.

\textbf{Principais conceitos.}
\begin{itemize}
    \item A amplificação de dados: poucos parâmetros geram grande variedade de conteúdo, inclusive sob demanda na GPU \cite[p.~1]{smelik2014}.
    \item A organização dos métodos por tipo de conteúdo (terreno, vegetação, água, estradas, cidades, edifícios, interiores) e a avaliação de cada um por facilidade de uso, controle e família de técnica \cite[p.~2]{smelik2014}.
    \item Uma seção sobre sistemas comerciais (Terragen, SpeedTree, CityEngine, Houdini) \cite[pp.~10--12]{smelik2014}.
    \item Os dois obstáculos à adoção: a falta de controle intuitivo e a falta de integração entre métodos \cite[pp.~12--15]{smelik2014}.
\end{itemize}

\textbf{Relação com o tema.} Mostra o procedural em escala de mundo e o problema de controle que também afeta as texturas procedurais. O levantamento exclui a geração de texturas \cite[p.~2]{smelik2014}, por isso é usado para o contexto de produção e para os produtos reais, e não para a texturização em si.

\textbf{Informações aproveitadas.} Amplificação de dados e custo da produção manual (quando usar), baixa controlabilidade e desempenho interativo (limitações e custo) (Cap.~\ref{chap:fundamentacao}); SpeedTree, CityEngine, Houdini e Substance Designer como produtos reais (Cap.~\ref{chap:aplicacoes}).

\section{Relação entre as publicações}

As três publicações cobrem camadas diferentes do problema. Heckbert trata de como uma coordenada \((u,v)\) chega a um pixel e quanto custa ler a textura. Dong et al.\ tratam do conteúdo procedural da textura em si. Smelik et al.\ tratam do procedural em escala de mundo e do problema de controle que limita seu uso. Nenhuma delas trata diretamente do UV dinâmico em tempo real. Esse ponto é coberto pela implementação do grupo e pelas fontes técnicas dos motores e estúdios (Capítulos~\ref{chap:vulkan}, \ref{chap:aplicacoes} e~\ref{chap:jogos}).
